\chapter{卷积、归一化与 Softmax 算子}

矩阵乘是算子开发的核心，但不是全部。推理引擎还需要卷积、归一化、Softmax、激活函数和残差加法等算子。它们的共同特点是：公式通常不难，性能和正确性细节很多。卷积考验数据布局和复用，归一化考验归约和内存流量，Softmax 考验数值稳定性。

\section{卷积的数据流}

卷积在视觉模型中常见。直接卷积会遍历 batch、输出通道、输出空间、输入通道和卷积核窗口。输入像素被相邻窗口复用，权重被多个输出位置复用。优化卷积时，必须决定数据布局、循环顺序和是否把卷积转化为 GEMM。

im2col 方法把每个卷积窗口展开成矩阵，再调用 GEMM。它用额外内存换取规则矩阵乘。直接卷积减少中间数据，但访问更复杂。Winograd 可以减少小卷积核乘法数量，但增加变换和数值误差。工程上没有永远最好的卷积算法，只有适合当前 shape、dtype 和硬件的选择。

\section{归一化}

LayerNorm 和 RMSNorm 是 Transformer 中常见算子。它们先对一段向量做归约，得到均值、方差或平方均值，再对每个元素缩放。归一化的 FLOPs 不一定高，但会扫描完整向量，常常受内存带宽和归约结构影响。

\begin{lstlisting}[language=C++,caption={RMSNorm 的教学实现}]
#include <cmath>
#include <cstdint>
#include <span>

void rms_norm(std::span<const float> x,
              std::span<const float> weight,
              std::span<float> y,
              float epsilon) {
    float square_sum = 0.0F;
    for (std::int64_t i = 0; i < static_cast<std::int64_t>(x.size()); ++i) {
        const float value = x[static_cast<std::size_t>(i)];
        square_sum += value * value;
    }
    const float mean_square = square_sum / static_cast<float>(x.size());
    const float scale = 1.0F / std::sqrt(mean_square + epsilon);
    for (std::int64_t i = 0; i < static_cast<std::int64_t>(x.size()); ++i) {
        y[static_cast<std::size_t>(i)] =
            x[static_cast<std::size_t>(i)] * scale *
            weight[static_cast<std::size_t>(i)];
    }
}
\end{lstlisting}

优化归一化时，要考虑多累加器、SIMD 归约、epsilon、低精度累加和融合。不要为了少一次扫描破坏数值稳定性。

\section{Softmax}

Softmax 把一组 logits 转成概率分布。稳定实现要先减最大值，再计算指数和归一化。直接计算 $\exp(x_i)$ 可能溢出；减去最大值不会改变概率比例，却能显著改善数值稳定性。

\begin{lstlisting}[language=C++,caption={稳定 Softmax 的基本结构}]
void softmax_stable(std::span<const float> input,
                    std::span<float> output) {
    float max_value = input[0];
    for (std::int64_t i = 1; i < static_cast<std::int64_t>(input.size()); ++i) {
        max_value = std::max(max_value, input[static_cast<std::size_t>(i)]);
    }
    float sum = 0.0F;
    for (std::int64_t i = 0; i < static_cast<std::int64_t>(input.size()); ++i) {
        const float value =
            std::exp(input[static_cast<std::size_t>(i)] - max_value);
        output[static_cast<std::size_t>(i)] = value;
        sum += value;
    }
    const float inv_sum = 1.0F / sum;
    for (std::int64_t i = 0; i < static_cast<std::int64_t>(output.size()); ++i) {
        output[static_cast<std::size_t>(i)] *= inv_sum;
    }
}
\end{lstlisting}

\section{和 LCQI 的关系}

LCQI 第一阶段不必实现卷积和完整 Attention，但要实现 ReLU、线性层和输出分类。后续扩展到 CNN 时加入卷积，扩展到语言模型时加入归一化、Softmax、RoPE 和 Attention。每加入一个算子，都要先写标量正确版本，再写测试，再考虑 SIMD、多线程和量化。

\begin{warningbox}
算子开发不要直接从最快版本开始。没有标量基线和误差检查，后续 SIMD 或多线程版本很难证明自己正确。
\end{warningbox}
